\documentclass[12pt]{article}
\usepackage[T1]{fontenc}
\usepackage[utf8]{inputenc}
\usepackage{lmodern}
\usepackage{textcomp}
\usepackage{enumitem}
\usepackage{geometry}
\geometry{margin=1in}
\begin{document}
\begin{center}
Answer Key - Geography - K-12 Level Assessment 19 \
\small (Answers are ordered exactly as in the question paper.)
\end{center}

\section*{Multiple Choice}
\begin{enumerate}[label=\arabic*.]
\item \textbf{Answer:} C
\\ \textit{Options:} A) Canada; B) United States; C) France; D) People's Republic of China
\\ \textit{Note:} France is not a federal state because it operates under a unitary system of government, where power is centralized in the national government rather than being divided between multiple levels of authority.
\item \textbf{Answer:} B
\\ \textit{Options:} A) Labor; B) Transport; C) Raw materials; D) Agglomeration costs
\\ \textit{Note:} Transport costs are the most important factor in locating an industry according to Alfred Weber's least-cost theory because they significantly influence overall production expenses and the ability to compete in the market by reducing the costs associated with moving raw materials to factories and finished goods to consumers.
\item \textbf{Answer:} C
\\ \textit{Options:} A) They are crossing another country's borders.; B) They are not taking anything with them.; C) They are moving permanently.; D) They will be gone more than one year.
\\ \textit{Note:} The correct answer is based on the key concept that migration refers to a permanent change of residence, which applies to Johnny and Susie's decision to move to Scotland to renovate and live in the family castle.
\item \textbf{Answer:} C
\\ \textit{Options:} A) Iran; B) Iraq; C) Turkey; D) Egypt
\\ \textit{Note:} The majority of Kurds are found in Turkey because approximately 15-20 million Kurds, which represents the largest concentration of the ethnic group, reside primarily in the southeastern part of the country.
\item \textbf{Answer:} C
\\ \textit{Options:} A) nationalism.; B) devolution.; C) supranationalism.; D) complementarity.
\\ \textit{Note:} Supranationalism refers to the process by which countries agree to work together and cede some of their sovereign powers to a higher authority in order to achieve common goals, making it the correct term for a voluntary association of nations.
\end{enumerate}

\section*{True / False}
\begin{enumerate}[label=\arabic*.]
\setcounter{enumi}{5}
\item \textbf{Answer:} True
\\ \textit{Options:} A) True; B) False
\\ \textit{Note:} Secularism is defined as the principle of separating religion from political, social, and educational institutions, leading to a societal framework that is either indifferent to religious considerations or actively promotes a non-religious perspective.
\item \textbf{Answer:} True
\\ \textit{Options:} A) True; B) False
\\ \textit{Note:} Urban agriculture primarily focuses on food production within city environments and does not inherently address or contribute to the purification or renewal of water supplies, which typically requires more specialized environmental management practices.
\item \textbf{Answer:} False
\\ \textit{Options:} A) True; B) False
\\ \textit{Note:} The First Agricultural Revolution primarily involved the domestication of plants and animals and the transition from nomadic lifestyles to settled farming practices, rather than the use of advanced technology and urban agriculture.
\item \textbf{Answer:} True
\\ \textit{Options:} A) True; B) False
\\ \textit{Note:} Canada's third major migration stream primarily consists of immigrants from Asia and Europe, rather than Latin America, which supports the truth of the statement.
\item \textbf{Answer:} True
\\ \textit{Options:} A) True; B) False
\\ \textit{Note:} A slash-and-burn agricultural society typically relies on small plots of land for farming, which can lead to soil depletion and limited resources over time, causing population pressure as the community grows and struggles to sustain itself.
\end{enumerate}

\section*{Long Form Response}
\begin{enumerate}[label=\arabic*.]
\setcounter{enumi}{10}
\item \textbf{Answer:} Bulk-gaining industries are those in which the final product weighs more or takes up more space than the raw materials used in its production. Bottled orange juice is a prime example of this type of industry, as the production process involves squeezing oranges to extract juice, which is then packaged in heavy glass or plastic bottles. The oranges themselves are lightweight compared to the bottled juice, making the final product bulkier. This process not only adds value to the oranges but also impacts the economy by creating jobs in agriculture, processing, and distribution. Additionally, the demand for bottled orange juice supports local farmers and contributes to the economy through exports.
\\ \textit{Note:} correct because it clearly illustrates how bottled orange juice exemplifies a bulk-gaining industry by demonstrating that the production process increases the weight and volume of the final product compared to the raw materials, thereby adding economic value.
\item \textbf{Answer:} China is industrializing rapidly, which has significantly impacted East Asia's economy, environment, and society. Economically, this swift industrialization has led to increased production and exports, making China a major player in global trade and boosting the economies of neighboring countries. However, the environmental consequences are severe, as industrial growth has resulted in pollution, deforestation, and climate change challenges. Socially, rapid industrialization has led to urban migration, changing lifestyles, and economic disparities, as many people move to cities for better job opportunities. Overall, while China's industrialization has spurred economic growth, it has also created pressing environmental and social issues that the region must address.
\\ \textit{Note:} correct because it effectively outlines the multifaceted impacts of China's rapid industrialization on East Asia, emphasizing economic growth and trade, environmental degradation, and social changes, thereby addressing the core elements of the question.
\end{enumerate}

\vfill
\noindent\textit{End of Answer Key}
\end{document}